\section{Truy vết}\label{ch:truyvet}

Đề bài yêu cầu truy vết từ yêu cầu đến user story hoặc use case, thiết kế, mã nguồn và test. Ma trận truy vết nằm trong \code{docs/\allowbreak{}traceability-matrix.\allowbreak{}md}, cập nhật đến \code{v1.1.1}. Ma trận được sửa cùng pull request với mã nguồn mỗi khi yêu cầu, thiết kế hoặc test thay đổi (Definition of Done, mục 5).

Ký hiệu thiết kế trong ma trận:
\begin{itemize}
\item \textbf{Activity-01:} activity diagram.
\item \textbf{Sequence-01:} sequence diagram.
\item \textbf{Class-01:} class diagram miền.
\item \textbf{ERD-01:} ERD.
\item \textbf{C4-01:} sơ đồ container.
\item \textbf{Deploy-01:} deployment diagram.
\item \textbf{ADR-00x:} quyết định kiến trúc.
\end{itemize}

Ký hiệu test theo bảng \ref{tab:tvkyhieu}.

\subsection{Yêu cầu chức năng}
{\scriptsize
\begin{longtable}{@{}L{2.4cm}L{2cm}L{2.6cm}L{3.6cm}L{3.3cm}@{}}
\caption{Ma trận truy vết yêu cầu chức năng}\label{tab:tvfr}\\
\toprule
\textbf{Mã yêu cầu} & \textbf{Use-case / User story} & \textbf{Thiết kế / Sơ đồ} & \textbf{Code / API} & \textbf{Test} \\ \midrule \endfirsthead
\toprule \textbf{Mã yêu cầu} & \textbf{Use-case / User story} & \textbf{Thiết kế / Sơ đồ} & \textbf{Code / API} & \textbf{Test} \\ \midrule \endhead
\bottomrule \endfoot
FR-01 Menu & US-01 Xem menu; UC-00 & Activity-01, Class-01 (Product) & \code{GET /products}; \code{apps/\allowbreak{}web/\allowbreak{}app/\allowbreak{}menu} & API-T03; UT \code{products.\allowbreak{}service.\allowbreak{}spec}; WEB \code{menu/\allowbreak{}page.\allowbreak{}test}; E2E-T01 \\
FR-02 Trạng thái menu & US-01 & Activity-01 & \code{apps/\allowbreak{}web/\allowbreak{}app/\allowbreak{}menu} (loading, rỗng, lỗi) & WEB \code{menu/\allowbreak{}page.\allowbreak{}test} \\
FR-03 Chi tiết, size, topping, đường, đá & US-02, US-06, US-23 & Activity-01, Class-01 (Product, Topping) & \code{GET /\allowbreak{}products/\allowbreak{}:id}; \code{apps/\allowbreak{}web/\allowbreak{}app/\allowbreak{}product/\allowbreak{}[id]} (tối đa 5 topping) & API-T04--T06; WEB \code{product/\allowbreak{}[id]/\allowbreak{}page.\allowbreak{}test}; E2E-T02 \\
FR-04 Giỏ hàng, tối đa 50 ly & US-06, US-23 & Activity-01 & \code{apps/\allowbreak{}web/\allowbreak{}app/\allowbreak{}cart}, \code{lib/cart.ts}; \code{POST /\allowbreak{}products/\allowbreak{}availability} & WEB \code{cart/\allowbreak{}page.\allowbreak{}test}; INT-R04; UT \code{availability.\allowbreak{}spec}; E2E-T02, E2E-T07 \\
FR-05 Đăng ký, đăng nhập, JWT, bcrypt & US-07; UC-01 (tiền điều kiện) & Sequence-01 (kiểm tra JWT); ADR-006 & \code{POST /\allowbreak{}auth/\allowbreak{}register}, \code{POST /\allowbreak{}auth/\allowbreak{}login}, \code{GET /auth/me}; \code{apps/\allowbreak{}web/\allowbreak{}app/\allowbreak{}login}, \code{register} & UT \code{auth.\allowbreak{}service.\allowbreak{}spec}; INT-A01--A02; AUTH-T02, T09; WEB \code{login/\allowbreak{}page.\allowbreak{}test}; E2E-T02 \\
FR-06 Chặn đặt đơn khi chưa đăng nhập & US-07; UC-01 (luồng ngoại lệ) & Activity-01 (nhánh chưa đăng nhập), Sequence-01 & \code{JwtAuthGuard}; checkout chuyển \code{/login?next=} & AUTH-T01; WEB \code{checkout/\allowbreak{}page.\allowbreak{}test}; E2E-T03 \\
FR-07 Tạo đơn PENDING sau khi kiểm tra nguyên liệu & US-08; UC-01 & Activity-01 (nhánh hết hàng), Sequence-01, Class-01 (Order, OrderItem) & \code{POST /orders}; \code{orders.\allowbreak{}service.\allowbreak{}ts}, \code{inventory/\allowbreak{}order-needs.\allowbreak{}ts} & INT-A03, A05, A06; INT-R06; AUTH-T06; UT \code{pricing.spec}, \code{order-needs.\allowbreak{}spec} \\
FR-08 Thanh toán Ví/Thẻ, thất bại cho thử lại & US-09; UC-01 & Sequence-01 (alt thành công/thất bại); ADR-004 & \code{POST /\allowbreak{}payments}; \code{payment-gateway.\allowbreak{}ts}; \code{apps/\allowbreak{}web/\allowbreak{}app/\allowbreak{}checkout} & UT \code{payment-gateway.\allowbreak{}spec}; INT-P01, P04; AUTH-T08; WEB \code{checkout/\allowbreak{}payment.\allowbreak{}test}, \code{recovery.test}; E2E-T02 \\
FR-09 Idempotency-Key & US-09; UC-01 & Sequence-01 (gửi lại yêu cầu); ADR-005 & \code{PaymentsService.\allowbreak{}settle}, \code{@@unique([userId, idempotencyKey])} & INT-P02, P03, P08, P09; INT-Q05, Q06; AUTH-T07 \\
FR-10 Máy trạng thái đơn & US-11, US-12, US-25 & Class-01 (OrderStatus) & \code{orders/\allowbreak{}order-status.\allowbreak{}ts}, \code{PATCH /\allowbreak{}orders/\allowbreak{}:id/\allowbreak{}status}; \code{apps/\allowbreak{}web/\allowbreak{}app/\allowbreak{}orders/\allowbreak{}[id]} & UT \code{order-status.\allowbreak{}spec}; INT-A10, P10; AUTH-T05; WEB \code{orders/\allowbreak{}[id]/\allowbreak{}page.\allowbreak{}test} \\
FR-11 Lịch sử, chi tiết đơn & US-11, US-25 & Class-01 & \code{GET /\allowbreak{}orders/\allowbreak{}me}, \code{GET /\allowbreak{}orders/\allowbreak{}:id}; \code{apps/\allowbreak{}web/\allowbreak{}app/\allowbreak{}orders/\allowbreak{}history}, \code{orders/[id]} & INT-A07, A08; WEB \code{orders/\allowbreak{}history/\allowbreak{}page.\allowbreak{}test}, \code{orders/\allowbreak{}[id]/\allowbreak{}page.\allowbreak{}test}; E2E-T02 \\
FR-12 Trừ kho có điều kiện, hoàn kho khi hủy & US-10 & Sequence-01 (transaction); ADR-005 & \code{inventory/\allowbreak{}stock.\allowbreak{}ts} (\code{UPDATE … WHERE stock >= q}) & INT-P01, P06, P07, P10; INT-Q07 \\
FR-13 Mã khuyến mãi & US-13 & Class-01 (Promotion) & \code{orders/\allowbreak{}pricing.\allowbreak{}ts}; ô mã ở checkout & UT \code{pricing.spec}; INT-A03, A04, A11 \\
FR-14 Điểm thưởng & US-14 & Class-01 (LoyaltyTransaction) & \code{loyalty/\allowbreak{}loyalty.\allowbreak{}ts} & INT-P01, P10; INT-Q04; INT-V02 \\
FR-15 Nhân viên xử lý đơn & US-12 & Class-01 & \code{GET /\allowbreak{}orders?status=}, \code{PATCH /\allowbreak{}orders/\allowbreak{}:id/\allowbreak{}status}; \code{apps/\allowbreak{}web/\allowbreak{}app/\allowbreak{}staff} (đăng xuất về trang chủ từ v1.1.1) & INT-A08--A10; AUTH-T03, T04; WEB \code{staff/\allowbreak{}staff.\allowbreak{}test} \\
FR-16 Định lượng theo công thức, size, topping & US-02, US-08, US-23, US-24 & ERD-01 (RecipeItem, ToppingRecipeItem) & \code{inventory/\allowbreak{}recipe-calculator.\allowbreak{}ts}, \code{availability.\allowbreak{}ts}; \code{POST /\allowbreak{}products/\allowbreak{}availability}; \code{GET /\allowbreak{}ingredients/\allowbreak{}recipes} & UT \code{recipe-calculator.\allowbreak{}spec}, \code{availability.\allowbreak{}spec}; INT-R01--R03, R08, R09 \\
FR-17 Kho nguyên liệu, nhập thêm & US-15, US-24 & ERD-01 (Ingredient) & \code{GET /\allowbreak{}ingredients}, \code{PATCH /\allowbreak{}ingredients/\allowbreak{}:id/\allowbreak{}restock}, \code{GET /\allowbreak{}ingredients/\allowbreak{}recipes}; \code{apps/\allowbreak{}web/\allowbreak{}app/\allowbreak{}staff/\allowbreak{}ingredients}, \code{staff/recipes} & UT \code{ingredient-status.\allowbreak{}spec}; INT-S01--S03; INT-R07, R09; WEB \code{staff/\allowbreak{}staff.\allowbreak{}test}; E2E-T05 \\
FR-18 Thanh toán QR mô phỏng & US-16 & ADR-004 & \code{POST /\allowbreak{}payments/\allowbreak{}qr}, \code{/\allowbreak{}payments/\allowbreak{}qr/\allowbreak{}:token…}; \code{apps/\allowbreak{}web/\allowbreak{}app/\allowbreak{}pay/\allowbreak{}qr/\allowbreak{}[token]}, \code{components/\allowbreak{}qr-payment-panel.\allowbreak{}tsx} & INT-Q01--Q09; WEB \code{pay/\allowbreak{}qr/\allowbreak{}[token]/\allowbreak{}page.\allowbreak{}test}, \code{qr-payment-panel.\allowbreak{}test}; E2E-T04 \\
FR-19 Đổi quà & US-17 & Class-01 (LoyaltyTransaction) & \code{GET /rewards}, \code{POST /\allowbreak{}rewards/\allowbreak{}:id/\allowbreak{}redeem}, \code{GET /\allowbreak{}rewards/\allowbreak{}me}; \code{apps/\allowbreak{}web/\allowbreak{}app/\allowbreak{}rewards} & UT \code{reward-code.\allowbreak{}spec}; INT-S04--S07; WEB \code{rewards/\allowbreak{}rewards.\allowbreak{}test} \\
FR-20 Tự hủy đơn quá hạn & US-18 & Class-01 (OrderStatus) & \code{orders/\allowbreak{}order-expiry.\allowbreak{}service.\allowbreak{}ts} & UT \code{order-expiry.\allowbreak{}service.\allowbreak{}spec}; INT-S08 \\
FR-21 VNPay sandbox & US-19 & ADR-004 & \code{payments/\allowbreak{}vnpay/\allowbreak{}*} (HMAC-SHA512, Return URL, IPN); \code{apps/\allowbreak{}web/\allowbreak{}app/\allowbreak{}checkout/\allowbreak{}vnpay-return}, \code{components/\allowbreak{}vnpay-button.\allowbreak{}tsx} & UT \code{vnpay-signature.\allowbreak{}spec}; INT-V01--V06; WEB \code{vnpay-return/\allowbreak{}page.\allowbreak{}test}, \code{vnpay-button.\allowbreak{}test}; kiểm thử thật trên production 16/11/2026 (bằng chứng task5 ảnh 16--26) \\
FR-22, FR-23 & Won't & -- & Không thực hiện & -- \\
\end{longtable}
}

\subsection{Yêu cầu phi chức năng}
{\scriptsize
\begin{longtable}{@{}L{2.4cm}L{1.4cm}L{2cm}L{3.8cm}L{4.3cm}@{}}
\caption{Ma trận truy vết yêu cầu phi chức năng}\label{tab:tvnfr}\\
\toprule
\textbf{Mã yêu cầu} & \textbf{User story} & \textbf{Thiết kế} & \textbf{Hiện thực} & \textbf{Kiểm chứng} \\ \midrule \endfirsthead
\toprule \textbf{Mã yêu cầu} & \textbf{User story} & \textbf{Thiết kế} & \textbf{Hiện thực} & \textbf{Kiểm chứng} \\ \midrule \endhead
\bottomrule \endfoot
NFR-01 Hiệu năng \code{GET /products} & US-01 & C4-01 & Một truy vấn có include, dữ liệu seed & Đo thủ công khi chạy Docker Compose (Sprint 1); chưa có test tự động \\
NFR-02 Băm mật khẩu, PIN; JWT có hạn; phân quyền & US-07 & ADR-006, Sequence-01 & bcrypt, \code{JWT\_\allowbreak{}EXPIRES\_\allowbreak{}IN}, \code{RolesGuard} & UT \code{auth.\allowbreak{}service.\allowbreak{}spec}; AUTH-T01--T05; INT-A01; INT-R07; INT-Q03 \\
NFR-03 Secret không nằm trong repo/log & US-21 & Deploy-01, ADR-007 & \code{.env.example} không có giá trị thật; GitHub Secrets, biến Render & Rà soát repo; \code{docs/\allowbreak{}deployment.\allowbreak{}md}; database production chặn kết nối Internet \\
NFR-04 Validate đầu vào hai phía & US-07, US-08 & ADR-002 & \code{ValidationPipe} (whitelist, forbidNonWhitelisted), DTO; kiểm tra form ở web & AUTH-T06, T08, T09; INT-R05; WEB \code{login/\allowbreak{}page.\allowbreak{}test}, \code{product/\allowbreak{}[id]/\allowbreak{}page.\allowbreak{}test} \\
NFR-05 Mobile-first, PWA & US-01, US-25 & C4-01 & Tailwind; \code{app/\allowbreak{}manifest.\allowbreak{}ts}, \code{app/icon.svg}, \code{public/\allowbreak{}icon-192.\allowbreak{}png}, \code{public/\allowbreak{}icon-512.\allowbreak{}png} & E2E-T06 (375px không cuộn ngang, manifest); WEB \code{page.test} (kích thước biểu tượng); bằng chứng task5 ảnh 62, 65 \\
NFR-06 Thông báo lỗi tiếng Việt & US-25 & ADR-008 & Thông báo trong service, \code{lib/api.ts} & Các test INT, WEB kiểm tra nội dung thông báo \\
NFR-07 \code{docker compose up} & US-21 & Deploy-01 (local) & \code{docker-compose.\allowbreak{}yml}, Dockerfile & Bằng chứng task5 ảnh 04 \\
NFR-08 CI/CD, rollback & US-04, US-21 & Deploy-01, ADR-007 & \code{.\allowbreak{}github/\allowbreak{}workflows/\allowbreak{}ci.\allowbreak{}yml} (2 job), \code{cd.yml} (CI $\rightarrow$ GHCR $\rightarrow$ Render $\rightarrow$ smoke test, rollback bằng \code{workflow\_\allowbreak{}dispatch}) & Lần chạy CI/CD trên GitHub Actions; bằng chứng task5 ảnh 02--03, 05, 14--15, 24, 27--28, 64 \\
NFR-09 \code{GET /health}, log & US-21 & ADR-008 & \code{health.\allowbreak{}controller.\allowbreak{}ts} (version, db) & UT \code{health.\allowbreak{}controller.\allowbreak{}spec}; API-T01, T02; smoke test CD \\
NFR-10 Migration, seed tái tạo được & US-03 & ERD-01, ADR-003 & \code{prisma/\allowbreak{}migrations}, \code{prisma/\allowbreak{}seed.\allowbreak{}ts} & CI migrate mỗi lần chạy; job E2E migrate và seed database mới \\
NFR-11 Lint, test tự động & US-04, US-20 & -- & ESLint, Vitest, Playwright & CI job "Install, Lint, Unit Test, Build" (bắt buộc) và "E2E (Playwright)" \\
\end{longtable}
}

\subsection{Use case UC-01 đến test}
Bảng \ref{tab:tvuc} đi theo từng bước của UC-01 và chỉ ra test kiểm chứng bước đó. Cả luồng chính lẫn luồng ngoại lệ quan trọng đều có ít nhất một integration test hoặc E2E.

{\small
\begin{longtable}{@{}L{5.4cm}L{2.6cm}L{5.9cm}@{}}
\caption{Truy vết các bước của UC-01 đến test}\label{tab:tvuc}\\
\toprule
\textbf{Bước của UC-01} & \textbf{Yêu cầu} & \textbf{Test} \\ \midrule \endfirsthead
\toprule \textbf{Bước của UC-01} & \textbf{Yêu cầu} & \textbf{Test} \\ \midrule \endhead
\bottomrule \endfoot
Chọn món, tùy chọn, thêm vào giỏ & FR-03, FR-04 & E2E-T02; WEB \code{product/\allowbreak{}[id]/\allowbreak{}page.\allowbreak{}test} \\
Chưa đăng nhập bị chuyển sang đăng nhập & FR-06 & E2E-T03; AUTH-T01 \\
Tạo đơn, kiểm tra nguyên liệu & FR-07, FR-16 & INT-A03, A05; INT-R06 \\
Thanh toán, idempotency & FR-08, FR-09 & INT-P01--P04, P08, P09; E2E-T02 \\
Thanh toán thất bại, thử lại & FR-08 & INT-P04; E2E-T02 \\
Trừ kho, cộng điểm & FR-12, FR-14 & INT-P01, P06, P07 \\
Xác nhận, theo dõi đơn & FR-10, FR-11 & E2E-T02; WEB \code{orders/\allowbreak{}[id]/\allowbreak{}page.\allowbreak{}test} \\
\end{longtable}
}

\subsection{Nhận xét về độ phủ}
\begin{itemize}
\item \textbf{Yêu cầu Must:} 11/12 FR Must có ít nhất một integration test hoặc E2E. FR-02 (trạng thái menu) là hành vi thuần giao diện nên được kiểm chứng bằng component test.
\item \textbf{NFR-01 (hiệu năng):} đến \code{v1.1.1} chỉ đo thủ công khi chạy Docker Compose ở Sprint 1. Sau đó nhóm bổ sung PERF-T01 chạy trong CI (mục \ref{sec:nfrtest}). Test này đo độ trễ tuần tự, chưa phải load test.
\item \textbf{NFR-03 (secret):} kiểm chứng bằng SEC-T01 (quét repository trong CI) kết hợp rà soát cấu hình GitHub Secrets và Render.
\item \textbf{NFR-07, NFR-08:} kiểm chứng bằng log chạy pipeline và ảnh bằng chứng.
\item \textbf{FR-22, FR-23 (Won't):} không có thiết kế và test, đúng phạm vi đã chốt.
\end{itemize}
